\documentclass[11pt]{article}
\usepackage[margin=1in]{geometry}
\usepackage{fontspec}
\usepackage{hyperref}
\usepackage{orcidlink}
\usepackage{fvextra}
\fvset{breaklines=true,breakanywhere=true,fontsize=\small}
\title{Dossier de Registro: padic-grn-search\\\large Búsqueda y certificación de ordenamientos p-ádicos}
\author{
Jesús Rogelio Pérez-Buendía\,\orcidlink{0000-0002-7739-4779}
\and
Víctor Nopal-Coello\,\orcidlink{0000-0003-2608-3636}
\and
Yuriria Cortés-Poza\,\orcidlink{0000-0003-1856-9206}
}
\date{2026-04-25}

\begin{document}
\maketitle
\begin{abstract}
\texttt{padic-grn-search} implementa estrategias de búsqueda para identificar y certificar ordenamientos relevantes en sistemas dinámicos discretos bajo el criterio de métricas p-ádicas. El dossier integra la justificación funcional del software, código ejecutable, pruebas automáticas y evidencia de corridas reproducibles para su registro formal.
\end{abstract}
\tableofcontents
\newpage

\section{Introducción}
Este software extiende el análisis del núcleo \texttt{padic-grn-core} mediante rutinas de optimización y exploración del espacio de permutaciones. Incluye métodos de búsqueda genética base y búsqueda exhaustiva para escenarios de tamaño controlado, priorizando verificabilidad de resultados y trazabilidad metodológica.

En términos de uso, \texttt{padic-grn-search} permite pasar de la evaluación de un ordering individual a la exploración sistemática de candidatos, facilitando estudios comparativos y procesos de validación en aplicaciones de redes de regulación genética y otros sistemas discretos.

\section{Datos generales}
\begin{itemize}
  \item Nombre: \texttt{padic-grn-search}
  \item Versión: \texttt{0.1.0}
  \item Lenguaje: Python 3.10+
  \item Tipo: optimización y certificación de orderings p-ádicos
  \item DOI: \texttt{DOI\_SOFTWARE\_SEARCH}
\end{itemize}

\subsection{Correos electrónicos oficiales de autores}
\begin{itemize}
  \item Jesús Rogelio Pérez-Buendía: \texttt{rogelio.perez@cimat.mx}
  \item Víctor Nopal-Coello: \texttt{victor\_nopal@uaeh.edu.mx}
  \item Yuriria Cortés-Poza: \texttt{yuriria.cortes@iimas.unam.mx}
\end{itemize}

\subsection{Financiamiento}
\begin{itemize}
  \item Modelos matemáticos y computacionales no convencionales para el estudio y análisis de problemas relevantes en Biología, CF 2019/217367.
\end{itemize}

\section{Autoría oficial}
\VerbatimInput{../../../AUTORIA-ORCID-DOI-OFICIAL.md}

\section{Documentación funcional completa}
\subsection{README}
\VerbatimInput{../README.md}
\subsection{Manual de usuario}
\VerbatimInput{../MANUAL-USUARIO.md}
\subsection{Manual técnico}
\VerbatimInput{../MANUAL-TECNICO.md}
\subsection{Configuración de paquete}
\VerbatimInput{../pyproject.toml}
\subsection{Cita académica}
\VerbatimInput{../CITATION.cff}
\subsection{Licencia}
\VerbatimInput{../LICENSE}
\subsection{Configuración de pruebas}
\VerbatimInput{../pytest.ini}

\section{Código fuente completo}
\subsection{Inicialización de paquete}
\VerbatimInput{../src/padic_grn_search/__init__.py}
\subsection{CLI}
\VerbatimInput{../src/padic_grn_search/cli.py}
\subsection{Motor de búsqueda}
\VerbatimInput{../src/padic_grn_search/engine.py}

\section{Pruebas completas}
\subsection{Prueba smoke}
\VerbatimInput{../tests/test_smoke.py}
\subsection{Pruebas de motor}
\VerbatimInput{../tests/test_engine.py}

\section{Evidencia textual de validación}
\subsection{Log de ejecución de pruebas}
\VerbatimInput{validacion_search_pytest.txt}
\subsection{Salida GA}
\VerbatimInput{../../../registro-dossier/06_evidencia_corridas/search_ga_p2.json}
\subsection{Salida exhaustiva min muE}
\VerbatimInput{../../../registro-dossier/06_evidencia_corridas/search_bbmin_p2.json}
\subsection{Salida exhaustiva max muA}
\VerbatimInput{../../../registro-dossier/06_evidencia_corridas/search_bbmax_p2.json}

\end{document}
